\documentclass[11pt,a4paper]{article}

\usepackage[utf8]{inputenc}
\usepackage{amsmath,amssymb,amsthm,bm,mathtools}
\usepackage{geometry}
\geometry{margin=1in}
\usepackage{booktabs}
\usepackage{array}
\usepackage{longtable}
\usepackage{url}
\usepackage{tikz}
\usetikzlibrary{calc,shapes.geometric,patterns,arrows.meta,positioning}
% mosaic_tikz_preamble.tex - Strict Two-Colour Symmetry Taxonomy for Composition Mosaics
%
% Colour Taxonomy:
%   - Symmetric / Reversible Compositions: mosaicblue (blue!30)
%   - Chiral / Non-Symmetric Compositions: mosaicpink (red!25)
%   - Opacity: fill opacity=0.65 for all polygon fills to enable interpenetration visibility.

\usepackage{tikz}
\usetikzlibrary{shapes.geometric, calc}

% Strict Two-Colour Taxonomy Definitions
\colorlet{mosaicblue}{blue!30}
\colorlet{mosaicpink}{red!25}

% TikZ Styles
\tikzset{
  % Standard uniform shell polygon size across all figures
  mosaic shell/.style={
    regular polygon,
    minimum size=1.8cm,
    draw=gray!40,
    thick
  },
  mosaic core shell/.style={
    regular polygon,
    minimum size=1.8cm,
    draw=gray!60,
    thick,
    densely dashed
  },
  mosaic polygon/.style={
    regular polygon,
    minimum size=1.8cm,
    draw=gray!40,
    thick
  },
  % Partition blocks
  symmetric block/.style={
    black, thick, fill=mosaicblue, fill opacity=0.65
  },
  chiral block/.style={
    black, thick, fill=mosaicpink, fill opacity=0.65
  },
  % Composition transversals (dashed paths)
  symm transversal/.style={
    blue!80!black,
    line width=1.5pt,
    dashed
  },
  symm dot/.style={
    fill=blue!80!black
  },
  chiral transversal/.style={
    red!75!black,
    line width=1.5pt,
    dashed
  },
  chiral dot/.style={
    fill=red!75!black
  },
  % Vertex coordinate discs and size badges
  active disc/.style={
    circle,
    fill=white,
    draw=black!70,
    line width=0.6pt,
    text=black,
    font=\tiny\bfseries,
    inner sep=1.0pt
  },
  pin dot/.style={
    circle,
    fill=white,
    draw=black!70,
    line width=0.6pt,
    text=black,
    font=\tiny\bfseries,
    inner sep=1.0pt
  },
  symm size badge/.style={
    circle,
    fill=mosaicblue,
    draw=blue!80!black,
    line width=0.6pt,
    text=blue!90!black,
    font=\tiny\bfseries,
    inner sep=0.7pt
  },
  chiral size badge/.style={
    circle,
    fill=mosaicpink,
    draw=red!75!black,
    line width=0.6pt,
    text=red!90!black,
    font=\tiny\bfseries,
    inner sep=0.7pt
  },
  coupling arrow/.style={
    ->,
    >=stealth,
    line width=1.1pt,
    shorten <=4pt,
    shorten >=4.5pt
  },
  inactive dot/.style={
    circle,
    fill=black!25,
    inner sep=1.0pt
  },
  mono size badge/.style={
    symm size badge
  },
  unigon dot/.style={
    black,
    fill=mosaicblue,
    line width=0.8pt
  },
  digon stroke/.style={
    line width=3.2pt,
    line cap=round
  },
  digon lens/.style={
    fill=mosaicblue,
    fill opacity=0.75,
    draw=blue!75!black,
    line width=0.8pt
  }
}

\usepackage{csquotes}
\MakeOuterQuote{"}
\usepackage{placeins}

\usepackage{hyperref}
\hypersetup{
  pdftitle={An Atlas of Hexadecagonal Compositions: Affine Macro-Orbits and Galois Multiplier Symmetries on Z16},
  pdfauthor={Gemma Knee and Antigra Vitti},
  colorlinks=true,
  linkcolor=blue!70!black,
  citecolor=blue!70!black,
  urlcolor=blue!70!black
}

\title{\textbf{An Atlas of Hexadecagonal Compositions:}\\ Affine Macro-Orbits and Galois Multiplier Symmetries on \texorpdfstring{$\mathbb{Z}_{16}$}{Z16}}
\author{Gemma Knee \and Antigra Vitti}
\date{November 2026}

% Custom Theorem Styles without terminating punctuation bacilli
\newtheoremstyle{breakdef}%
  {9pt}% Space above
  {9pt}% Space below
  {\normalfont}% Body font
  {0pt}% Indent amount
  {\bfseries}% Theorem head font
  {}% Punctuation after theorem head
  {\newline}% Space after theorem head
  {\thmname{#1}\thmnumber{ #2}\thmnote{ (#3)}}% Head spec

\newtheoremstyle{breakthm}%
  {9pt}% Space above
  {9pt}% Space below
  {\itshape}% Body font
  {0pt}% Indent amount
  {\bfseries}% Theorem head font
  {}% Punctuation after theorem head
  {\newline}% Space after theorem head
  {\thmname{#1}\thmnumber{ #2}\thmnote{ (#3)}}% Head spec

\theoremstyle{breakthm}
\newtheorem{theorem}{Theorem}[section]
\newtheorem{proposition}[theorem]{Proposition}
\newtheorem{corollary}[theorem]{Corollary}
\newtheorem{lemma}[theorem]{Lemma}

\theoremstyle{breakdef}
\newtheorem{definition}[theorem]{Definition}
\newtheorem{example}[theorem]{Example}
\newtheorem{remark}[theorem]{Remark}

\numberwithin{equation}{section}

\begin{document}

\maketitle

\begin{abstract}
On the 16-point cyclic lattice $\mathbb{Z}_{16}$, the space of equal-block set partitions of shape $z_4^4$ comprises $2{,}627{,}625$ raw mosaic partitions and $83{,}488$ non-equivalent dihedral orbits under $D_{16}$. While full graphical cataloguing of 83 thousand individual orbits would require over a thousand pages, the action of the affine group $\operatorname{Aff}(\mathbb{Z}_{16}) = \mathbb{Z}_{16} \rtimes (\mathbb{Z}_{16})^\times$ induces an outer automorphism group $G_{\text{macro}} = \operatorname{Aff}(\mathbb{Z}_{16}) / D_{16} \cong (\mathbb{Z}_{16})^\times / \{\pm 1\} \cong C_4$. This cyclic group of order 4, generated by the Galois dilation multiplier $\mu_3 \colon x \mapsto 3x \pmod{16}$, partitions the entire space into exactly $21{,}030$ Affine Macro-Orbits. We prove the exact spectral decomposition: $14$ monomers (size 1 fixed points), $295$ dimers (size 2 cycles), and $20{,}721$ tetramers (size 4 cycles). Cataloguing one canonical exemplar per macro-orbit collapses the visual universe to approximately $220$ pages. This monograph provides the complete mathematical foundations, the cataloguing conventions, and the inventory of singular symmetries.
\end{abstract}

\tableofcontents
\newpage

\section{Hexadecagonal Architecture}
\label{sec:introduction}

The sixteen vertices of the cyclic lattice $\mathbb{Z}_{16}$ give rise to an immense combinatorial landscape. When we partition the regular hexadecagon into four disjoint blocks of four vertices each (the shape $z_4^4$), the vertices connect to form inscribed polygons whose cyclic edge strides define integer compositions of $16$~\cite{Fripertinger1997,Kerber1999}. 

Across the full combinatorial universe, exactly $2{,}627{,}625$ distinct set partitions tile the sixteen vertices into four quadrilaterals:
\begin{equation}
\label{eqn:raw-universe-z44}
|\mathcal{P}(z_4^4)| = \frac{1}{4!} \binom{16}{4} \binom{12}{4} \binom{8}{4} \binom{4}{4} = 2{,}627{,}625
\end{equation}
Under the physical symmetries of the polygon—the translations and reflections generating the dihedral group $D_{16}$ of order $32$—these configurations gather into the graded bivariate orbit polynomial:
\begin{equation}
\label{eqn:orbit-poly-d16}
P_{16}(z_4^4, \omega) = 80802\omega^{32} + 2568\omega^{16} + 102\omega^{8} + 13\omega^{4} + 2\omega^{2} + \omega^{1}
\end{equation}
Evaluating this polynomial at $\omega = 1$ reveals the total census of non-equivalent dihedral mosaic orbits:
\begin{equation}
\label{eqn:total-d16-orbits}
P_{16}(z_4^4, 1) = 80802 + 2568 + 102 + 13 + 2 + 1 = 83{,}488
\end{equation}

Beneath this vast collection of 83 thousand geometric configurations lies a deeper, unifying algebra governed by modular multiplication. The affine group $\operatorname{Aff}(\mathbb{Z}_{16}) = \mathbb{Z}_{16} \rtimes (\mathbb{Z}_{16})^\times$ acts on the polygon not merely by turning or flipping it, but by dilating its chords through modular unit multipliers.

The outer automorphism group $G_{\text{macro}} = \operatorname{Aff}(\mathbb{Z}_{16}) / D_{16} \cong C_4$ acts upon the dihedral orbits as a cyclic permutation group of order $4$, driven by the Galois multiplier $\mu_3 \colon x \mapsto 3x \pmod{16}$. This multiplicative symmetry weaves the $83{,}488$ individual dihedral mosaics into exactly $21{,}030$ Affine Macro-Orbits.

\subsection{Lattice Geometry and Clock-Face Orientation}
\label{sec:clock-face-orientation}

Across all diagrams, plates, and catalogue listings in this atlas, the sixteen hexadecagonal vertices are identified as $\{0, 1, \dots, 15\}$ according to a clock-face geometry with $0$ at the top, advancing by $1$ clockwise.

A representative partition of the set of vertices as a \textit{mosaical shape} $z_4^4$ is exemplified by $P = \{\{0, 2, 7, 9\}, \{1, 4, 10, 11\}, \{3, 8, 12, 15\}, \{5, 6, 13, 14\}\}$.

A perimeter \textit{stride} between successive vertices of an inscribed block $B = \{v_1, v_2, v_3, v_4\}$ (ordered circularly $v_1 < v_2 < v_3 < v_4$) is defined as the forward modular step around the polygon, $s_i = (v_{i+1} - v_i) \bmod 16$ (where $v_5 \equiv v_1$) such that $\sum_{i=1}^4 s_i = 16$. For example, the first of the four blocks of our exemplar $P$, above, has strides $(2 - 0)$, $(7 - 2)$, $(9 - 7)$, $(0 - 9)$---which is to say $2, 5, 2, 7$ since $7 \equiv -9 \pmod{16}$---notice that $2 + 5 + 2 + 7 = 16$ as required.

Figure~\ref{fig:clock-face-convention} below illustrates vertex numbering, the entire example partition $P$, the vertices of the first block of $P$, and finally the strides of that first block.

\begin{figure}[htbp]
\centering
\input{tkz/T00646.tkz}
\caption{\label{fig:clock-face-convention}$z_4^4$-mosaic $P$ on $\mathbb{Z}_{16}$, with its modular arithmetic}
\end{figure}

The other three blocks of $P$ have corresponding strides $3 + 6 + 1 + 6$, $5 + 4 + 3 + 4$, and $1 + 7 + 1 + 7$.

Mosaics will generally be abbreviated (to save space) using hexadecimal numerals where $\mathrm{A}=10$, $\mathrm{B}=11$, $\mathrm{C}=12$, etc. For instance, the mosaic $P$ above is notated as \path{/0279/14AB/38CF/56DE/}, using lexicographical ordering to determine a \textit{canonical representative} of dihedrally equivalent partitions (mosaics).

\section{Affine Macro-Orbit Theory on \texorpdfstring{$\mathbb{Z}_{16}$}{Z16}}
\label{sec:affine-macro-orbits}

\subsection{The Outer Galois Multiplier Group}

Let $\operatorname{Aff}(\mathbb{Z}_n) = \{x \mapsto ax + b \pmod n \mid a \in (\mathbb{Z}_n)^\times, b \in \mathbb{Z}_n\}$ denote the affine group of the cyclic ring $\mathbb{Z}_n$. The dihedral group $D_n = \{x \mapsto \pm x + b \pmod n\}$ is a normal subgroup of $\operatorname{Aff}(\mathbb{Z}_n)$. Consequently, the outer quotient group:
\begin{equation}
\label{eqn:g-macro-def}
G_{\text{macro}} = \operatorname{Aff}(\mathbb{Z}_n) / D_n \cong (\mathbb{Z}_n)^\times / \{\pm 1\}
\end{equation}
acts naturally as a group of permutations on the set of $D_n$-orbits of composition mosaics.

\begin{proposition}[Cyclic Structure of the Macro-Group on $\mathbb{Z}_{16}$]
\label{prop:g-macro-c4}
For $n = 16$, the unit group $(\mathbb{Z}_{16})^\times = \{1, 3, 5, 7, 9, 11, 13, 15\}$ is isomorphic to $C_2 \times C_4$. Factoring out the reflection subgroup $\{\pm 1\} = \{1, 15\}$, the quotient group is cyclic of order $4$:
\begin{equation}
\label{eqn:c4-quotient}
G_{\text{macro}} \cong C_4 = \langle \mu_3 \rangle
\end{equation}
generated by the Galois dilation multiplier:
\begin{equation}
\label{eqn:mu3-def}
\mu_3 \colon \mathbb{Z}_{16} \to \mathbb{Z}_{16}, \quad x \mapsto 3x \pmod{16}
\end{equation}
\end{proposition}

\begin{proof}
The cosets of $(\mathbb{Z}_{16})^\times$ modulo $\{1, 15\}$ are:
\begin{enumerate}
    \item $C_0 = \{1, 15\}$ (identity)
    \item $C_1 = \{3, 13\}$
    \item $C_2 = \{7, 9\}$
    \item $C_3 = \{5, 11\}$
\end{enumerate}
Under multiplication modulo $16$, $3^1 \equiv 3 \in C_1$, $3^2 = 9 \equiv -7 \in C_2$, $3^3 = 27 \equiv 11 \equiv -5 \in C_3$, and $3^4 = 81 \equiv 1 \in C_0$. Thus $C_1$ has order $4$, generating the cyclic quotient $C_4$.
\end{proof}

The unique order-2 subgroup of $G_{\text{macro}}$ is generated by the involution:
\begin{equation}
\label{eqn:mu7-involution}
\mu_7 \colon x \mapsto 7x \equiv -9x \pmod{16}
\end{equation}
which corresponds to the intermediate affine subgroup:
\begin{equation}
\label{eqn:h2-subgroup}
H_2 = D_{16} \rtimes \langle \mu_7 \rangle \le \operatorname{Aff}(\mathbb{Z}_{16})
\end{equation}
of order $64$, whose unit set is $\{1, 7, 9, 15\}$.

We demonstrate such an $x \mapsto 3x \pmod{16}$ inflation with the first block in our example mosaic, $P$, above. This block is the subset $\{0, 2, 7, 9\}$ and is dilated (or inflated) by multiplying each vertex number by $3$, yielding $\{0, 6, 21, 27\}$. The last two elements---having been inflated round the clock, as it were---are not in our ground set except as vertices $5$ and $11$ (via modular arithmetic). So the dilated shape is now the subset $\{0, 5, 6, 11\}$---clearly not the same shape. Dilating again produces $0, 15, 18, 33$---the subset $\{0, 1, 2, 15\}$---and a third dilation yields $0, 3, 6, 45$, i.e.\ $\{0, 3, 6, 13\}$. We know, from the proof above, that a fourth dilation (having come this far without return) \textit{must} return us to the beginning, and $0, 9, 18, 39$---being (modulo $16$) the set $\{0, 2, 7, 9\}$---is indeed that first subset.

Figure~\ref{fig:inflation-in-action} below illustrates the four dilations of the first subset of partition $P$.

\begin{figure}[htbp]
\centering
\input{tkz/T00646-1-dil.tkz}
\caption{\label{fig:inflation-in-action}Successive inflations of subset $\{0, 2, 7, 9\}$}
\end{figure}

The second block of $P$ is $\{1, 4, 10, 11\}$, which maps to numbers $3, 12, 30, 33$---the subset $\{1, 3, 12, 14\}$---which maps again to $3, 9, 36, 42$---the subset $\{3, 4, 9, 10\}$. This maps, for the third time, to $9, 12, 27, 30$---the subset $\{9, 11, 12, 14\}$---which maps for the fourth (and perforce last) time to $27, 33, 36, 42$---the subset $\{1, 4, 10, 11\}$, which is the inevitable second block of $P$ again. The reader should verify that the remaining two blocks, i.e.\ $\{3, 8, 12, 15\}$ and $\{5, 6, 13, 14\}$, perform the same way, and also that at no point in any of the four dilations do the subsets intersect, meaning each stage produces a genuine partition of $\mathbb{Z}_{16}$.

\subsection{Burnside Spectral Decomposition}

Because $G_{\text{macro}} \cong C_4$ acts on the $83{,}488$ dihedral orbits, orbit lengths must divide $|C_4| = 4$. Every affine macro-orbit is therefore one of three types:
\begin{enumerate}
    \item \textbf{Monomers (Size 1)}: Fixed points invariant under all multipliers $\mu_3, \mu_5, \mu_7$.
    \item \textbf{Dimers (Size 2)}: 2-cycles invariant under the involution $\mu_7$, but transposed by $\mu_3 \leftrightarrow \mu_5$.
    \item \textbf{Tetramers (Size 4)}: Full 4-cycles under $\mu_3$:
    \[
    O \xrightarrow{\mu_3} \mu_3(O) \xrightarrow{\mu_3} \mu_7(O) \xrightarrow{\mu_3} \mu_5(O) \xrightarrow{\mu_3} O
    \]
\end{enumerate}

\begin{theorem}[Spectral Decomposition Theorem for $z_4^4$]
\label{thm:spectral-decomposition-z44}
Under the action of $G_{\text{macro}} \cong C_4$, the $83{,}488$ dihedral orbits of shape $z_4^4$ on $\mathbb{Z}_{16}$ decompose into exactly $21{,}030$ Affine Macro-Orbits, comprising:
\begin{itemize}
    \item $c_1 = 14$ monomers (size 1 fixed points)
    \item $c_2 = 295$ dimers (size 2 cycles)
    \item $c_4 = 20{,}721$ tetramers (size 4 cycles)
\end{itemize}
satisfying:
\begin{equation}
\label{eqn:spectral-sum}
c_1 \times 1 + c_2 \times 2 + c_4 \times 4 = 14 + 590 + 82{,}884 = 83{,}488
\end{equation}
and:
\begin{equation}
\label{eqn:total-macro-sum}
c_1 + c_2 + c_4 = 14 + 295 + 20{,}721 = 21{,}030
\end{equation}
\end{theorem}

\begin{proof}
Let $X$ denote the set of $2{,}627{,}625$ raw partitions of $\mathbb{Z}_{16}$ of shape $z_4^4$. The group $\operatorname{Aff}(\mathbb{Z}_{16})$ acts on $X$. The $128$ group elements partition into $9$ cycle classes in the symmetric group $S_{16}$. For each permutation $g \in \operatorname{Aff}(\mathbb{Z}_{16})$, the number of invariant set partitions $|\operatorname{Fix}_X(g)|$ is determined by cycle compatibility between $g$ and the block permutations in $S_4$:
\begin{center}
\begin{tabular}{llrr}
\toprule
Cycle Class of $g \in S_{16}$ & Elements in $\operatorname{Aff}_{16}$ & Fixed Partitions $|\operatorname{Fix}_X(g)|$ \\
\midrule
$(1^{16})$ & $1$ (identity) & $2{,}627{,}625$ \\
$(1^8, 2^4)$ & $2$ & $5{,}425$ \\
$(2^8)$ & $9$ & $4{,}025$ \\
$(1^2, 2^7)$ & $16$ & $945$ \\
$(4^4)$ & $12$ & $89$ \\
$(1^2, 2^3, 4^2)$ & $16$ & $33$ \\
$(1^4, 2^2, 4^2)$ & $8$ & $29$ \\
$(8^2)$ & $32$ & $5$ \\
$(16^1)$ & $32$ & $1$ \\
\bottomrule
\end{tabular}
\end{center}
Applying Burnside's Lemma~\cite{Burnside1911}:
\begin{enumerate}
    \item Dihedral group $D_{16}$ (order $32$, units $\{1, 15\}$):
    \[
    N_{D_{16}} = \frac{1}{32} \sum_{g \in D_{16}} |\operatorname{Fix}_X(g)| = 83{,}488
    \]
    \item Intermediate subgroup $H_2$ (order $64$, units $\{1, 7, 9, 15\}$):
    \[
    N_{H_2} = \frac{1}{64} \sum_{g \in H_2} |\operatorname{Fix}_X(g)| = 42{,}046
    \]
    \item Full affine group $\operatorname{Aff}_{16}$ (order $128$, all $8$ units):
    \[
    N_{\operatorname{Aff}_{16}} = \frac{1}{128} \sum_{g \in \operatorname{Aff}_{16}} |\operatorname{Fix}_X(g)| = 21{,}030
    \]
\end{enumerate}
The orbit counts under $D_{16}$, $H_2$, and $\operatorname{Aff}_{16}$ satisfy the linear system:
\[
\begin{cases}
c_1 + 2c_2 + 4c_4 = 83{,}488 \\
c_1 + 2c_2 + 2c_4 = 42{,}046 \\
c_1 + c_2 + c_4 = 21{,}030
\end{cases}
\]
Subtracting the second equation from the first yields $2c_4 = 83{,}488 - 42{,}046 = 41{,}442$, so $c_4 = 20{,}721$. Subtracting the third equation from the second yields $c_2 + c_4 = 42{,}046 - 21{,}030 = 21{,}016$, giving $c_2 = 21{,}016 - 20{,}721 = 295$. Finally, $c_1 = 21{,}030 - 295 - 20{,}721 = 14$.
\end{proof}

\section{Singular Symmetries: Monomers and Dimers}
\label{sec:singular-symmetries}

\subsection{The 14 Monomers (Fixed-Point Macro-Orbits)}

A dihedral orbit $O$ is an affine fixed point (monomer) if it is invariant under the entire outer automorphism group $G_{\text{macro}} \cong C_4$. For any partition $P \in O$, every affine multiplier $\mu_a$ maps $P$ to another member of the same dihedral orbit: $\mu_a(P) \in O$. 

By Theorem~\ref{thm:spectral-decomposition-z44}, there are exactly $14$ such fixed-point orbits. Their representative mosaics are exhibited in Figure~\ref{fig:14-monomers}, each badged directly with its dihedral orbit size $|O|$ ($1$, $4$, $16$, or $32$). These fourteen patterns form the absolute singular bedrock of the $z_4^4$ universe.

\begin{figure}[htbp]
\centering
\input{tkz/fig_14_monomers.tkz}
\caption{\label{fig:14-monomers}The $14$ affine monomer (fixed-point) dihedral orbits of shape $z_4^4$ on $\mathbb{Z}_{16}$, showing their respective dihedral orbit sizes}
\end{figure}

The $14$ monomers fall into four distinct dihedral size classes:
\begin{enumerate}
    \item \textbf{Size 1 ($M_1$)}:
    The unique regular coset partition:
    \[
    P_0 = \{\{0, 4, 8, 12\}, \{1, 5, 9, 13\}, \{2, 6, 10, 14\}, \{3, 7, 11, 15\}\}
    \]
    Each block is a regular square with stride profile $(4, 4, 4, 4)$ (shaded blue). It is invariant under all $32$ elements of $D_{16}$ and all $128$ elements of $\operatorname{Aff}(\mathbb{Z}_{16})$.
    
    \item \textbf{Size 4 ($M_2, M_3, M_4$)}:
    Three antipodal rectangular and rhombic configurations with dihedral stabiliser of order $8$:
    \begin{itemize}
        \item $M_2$: pairs of opposite rectangular blocks with stride profiles $(1, 7, 1, 7)$ and $(3, 5, 3, 5)$.
        \item $M_3$: four congruent rhombic blocks with uniform stride $(2, 6, 2, 6)$.
        \item $M_4$: hybrid pairing consisting of two regular squares $(4, 4, 4, 4)$ and two rhombi $(2, 6, 2, 6)$.
    \end{itemize}
    
    \item \textbf{Size 16 ($M_5, M_6, M_7, M_8$)}:
    Four orbits with dihedral stabiliser of order $2$:
    \begin{itemize}
        \item $M_5, M_6, M_7$: bilateral reflectional symmetry, each possessing a single reflection axis and no non-trivial rotations.
        \item $M_8$: central antipodal inversion symmetry, invariant under half-turn rotation $x \mapsto x + 8$ with no reflectional axes.
    \end{itemize}
    
    \item \textbf{Size 32 ($M_9 \dots M_{14}$)}:\newline
    Six fully asymmetric dihedral orbits with trivial dihedral stabiliser $|\operatorname{Stab}_{D_{16}}(P)| = 1$. Although entirely devoid of non-trivial dihedral symmetries, their stabiliser in the ambient affine group $\operatorname{Aff}(\mathbb{Z}_{16})$ has order $4$, generating an outer cyclic symmetry under $\mu_3$.
\end{enumerate}

\subsection{The Stabiliser Spectrum: The 1, 3, 4, 6 Decomposition}

The arithmetic decomposition of the fourteen monomers into dihedral orbit sizes:
\[
14 = 1 + 3[\times 4] + 4[\times 16] + 6[\times 32]
\]
is a direct consequence of the subgroup lattice of $D_{16}$ and the compatibility requirements imposed by the affine multiplier $\mu_3 \colon x \mapsto 3x \pmod{16}$.

More fundamentally, when indexed systematically across all six divisors of $|D_{16}| = 32$ in increasing order of dihedral orbit size $d \in \{1, 2, 4, 8, 16, 32\}$ (or decreasing stabiliser order $32/d \in \{32, 16, 8, 4, 2, 1\}$), the full coefficient vector is:
\[
(1, 0, 3, 0, 4, 6)
\]
In Fripertinger's cycle index framework~\cite{Fripertinger1997,Kerber1999}, these values appear directly as the $\omega$-coefficients of the shape monomial $z^\lambda = z_4^4$ under the Burnside-Polya plethysm, where the vanished coefficients at divisors $2$ and $8$ reflect the algebraic obstruction to isolating $\langle r^2 \rangle$ and $\langle r^4 \rangle$ symmetries under Galois dilation.

By the Orbit-Stabiliser Theorem, the size of a dihedral orbit $|O|$ and the order of the stabiliser subgroup $\operatorname{Stab}_{D_{16}}(P)$ of any representative partition $P \in O$ satisfy:
\[
|O| \cdot |\operatorname{Stab}_{D_{16}}(P)| = |D_{16}| = 32
\]
The possible subgroup orders in $D_{16}$ are divisors of $32$: $1, 2, 4, 8, 16, 32$. However, among the affine monomers, only stabiliser orders $32, 8, 2, 1$ occur. Stabiliser orders $16$ (orbit size $2$) and $4$ (orbit size $8$) are completely absent:

\begin{itemize}
    \item \textbf{Order 32 ($|O| = 1$, 1 monomer)}:
    The full dihedral group $D_{16}$ requires invariance under the generator of rotations $r \colon x \mapsto x + 1$. Because $r$ acts transitively on $\mathbb{Z}_{16}$, the four blocks of size $4$ must form an orbit under $r$, forcing each individual block to be invariant under $r^4 \colon x \mapsto x + 4$. The only subsets of size $4$ invariant under translation by $4$ are the cosets of the principal ideal $4\mathbb{Z}_{16}$. This configuration is unique, giving exactly $1$ monomer ($M_1$).
    
    \item \textbf{Order 8 ($|O| = 4$, 3 monomers)}:
    A stabiliser of order $8$ containing four rotations must have rotational subgroup $\langle r^4 \rangle \cong C_4$. Translation by $4$ permutes the four blocks as an element of the symmetric group $S_4$. Since $(r^4)^2 = r^8 \neq \operatorname{id}$, the permutation cannot be identity on all blocks without reducing to $M_1$. The permissible cycle structures in $S_4$ are:
    \begin{enumerate}
        \item Two 2-cycles $(B_a B_b)(B_c B_d)$: two distinct geometric arrangements arise, depending on whether the antipodal pairs are adjacent chords ($M_2$, strides $1$ and $7$) or alternate chords ($M_3$, strides $2$ and $6$).
        \item One 2-cycle and two fixed blocks $(B_a B_b)(B_c)(B_d)$: exactly one hybrid arrangement is realisable ($M_4$), where two blocks are invariant under $r^4$ (squares) and two blocks are interchanged (rhombi).
    \end{enumerate}
    Together, these yield precisely $2 + 1 = 3$ monomers.
    
    \item \textbf{Absence of Order 16 ($|O| = 2$) and Order 4 ($|O| = 8$)}:
    An orbit of size $2$ would require a stabiliser of order $16$, which contains the rotation subgroup $\langle r^2 \rangle \cong C_8$. Any partition invariant under translation by $2$ has its blocks composed of pairs $\{x, x+8\}$; under the action of the affine multiplier $\mu_3$, this translation group automatically forces full invariance under $r^4$ and reflections, collapsing the orbit to size $4$ or $1$. 
    
    Similarly, an order-$4$ stabiliser in $D_{16}$ would be either cyclic $\langle r^4 \rangle$ without reflections, or $\langle r^8, s \rangle$. However, the affine equivariance condition $\mu_3 \operatorname{Stab}_{D_{16}}(P) \mu_3^{-1} = \operatorname{Stab}_{D_{16}}(P)$ forces any partition possessing $\langle r^4 \rangle$ symmetry to also admit four reflection axes, lifting the stabiliser order from $4$ to $8$.
    
    \item \textbf{Order 2 ($|O| = 16$, 4 monomers)}:
    A stabiliser of order $2$ in $D_{16}$ is generated either by a single reflection $s$ ($\det = -1$) or by the central half-turn $r^8 \colon x \mapsto x + 8$ ($\det = +1$). Exactly $3$ monomers ($M_5, M_6, M_7$) possess a unique reflection axis, while exactly $1$ monomer ($M_8$) possesses central point-inversion symmetry $\langle r^8 \rangle$. In total, $3 + 1 = 4$ monomers emerge.
    
    \item \textbf{Order 1 ($|O| = 32$, 6 monomers)}:
    The remaining $6$ monomers have no non-trivial dihedral symmetries whatsoever. Each partition is completely chiral and asymmetric under rigid Euclidean motions. Nevertheless, in the ambient affine group $\operatorname{Aff}(\mathbb{Z}_{16})$, each possesses a cyclic stabiliser of order $4$ generated by an affine spiral $x \mapsto 3x + b$ or $x \mapsto 5x + b$.
\end{itemize}

\subsection{The 295 Dimers (2-Cycles)}

A dimer is an affine macro-orbit consisting of a pair of dihedral orbits $\{O_a, O_b\}$ that are interchanged by $\mu_3$ and $\mu_5$, but individually fixed by the involution $\mu_7 \colon x \mapsto 7x \equiv -9x \pmod{16}$:
\begin{equation}
\label{eqn:dimer-action}
\mu_3(O_a) = O_b, \quad \mu_3(O_b) = O_a, \quad \mu_7(O_a) = O_a, \quad \mu_7(O_b) = O_b
\end{equation}
Theorem~\ref{thm:spectral-decomposition-z44} proves that exactly $295$ such pairs exist, accounting for $590$ dihedral orbits ($0.7\%$ of the total). Because they are invariant under the order-64 group $H_2 = D_{16} \rtimes \langle \mu_7 \rangle$, their block composition profiles are preserved under the chord dilation $d \mapsto 7d \pmod{16}$.

\section{The Four-Cycle Tetramers and Comparative Scale Models}
\label{sec:tetramers-and-scale-models}

The vast majority of composition mosaics on $\mathbb{Z}_{16}$ (specifically, $82{,}884$ of the $83{,}488$ dihedral orbits, or $99.28\%$) gather into closed harmonic four-cycles under the Galois multiplier $\mu_3 \colon x \mapsto 3x \pmod{16}$. Theorem~\ref{thm:spectral-decomposition-z44} establishes that exactly $20{,}721$ such affine tetramers exist.

Within each tetramer cartouche, a complete orbit unfolds in a closed four-cycle:
\begin{equation}
\label{eqn:tetramer-progression}
\mu_1(O) \quad\longrightarrow\quad \mu_3(O) \quad\longrightarrow\quad \mu_7(O) \quad\longrightarrow\quad \mu_5(O) \quad\longrightarrow\quad \mu_1(O)
\end{equation}
Every step corresponds to multiplication by $3 \pmod{16}$, closing after four iterations because $3^4 \equiv 81 \equiv 1 \pmod{16}$.

\subsection{The Horizontal Linear Cartouche Geometry}

To make the algebraic sequence immediately and unambiguously readable, each tetramer unfolds in a horizontal linear cartouche ($1 \times 4$ strip). The four mosaics progress naturally from left to right across the page:
\begin{center}
\small
\begin{tabular}{cccc}
\toprule
\textbf{Position 1} & \textbf{Position 2} & \textbf{Position 3} & \textbf{Position 4} \\
\midrule
Root Orbit & First Dilation & Second Dilation (Involution) & Third Dilation (Inverse) \\
$\mu_1 \colon x \mapsto x$ & $\mu_3 \colon x \mapsto 3x$ & $\mu_7 \colon x \mapsto 7x$ & $\mu_5 \colon x \mapsto 5x$ \\
\bottomrule
\end{tabular}
\end{center}
Reading from left to right traces the exact dilation sequence:
\[
\mu_1(O) \;\xrightarrow{\mu_3}\; \mu_3(O) \;\xrightarrow{\mu_3}\; \mu_7(O) \;\xrightarrow{\mu_3}\; \mu_5(O) \;\xrightarrow{\mu_3}\; \mu_1(O)
\]
This linear orientation completely eliminates the ambiguity of two-dimensional grid layouts (where reading order conflicts with cyclic flow), making the location of the canonical root mosaic ($\mu_1$) and each subsequent modular dilation immediately discernible to the eye.

Following the clock-face geometry established in Section~\ref{sec:clock-face-orientation} (with vertex $0$ at 12~o'clock and indices advancing clockwise), each inscribed quadrilateral is strictly convex, with its vertices circularly ordered ($v_1 < v_2 < v_3 < v_4$) to avoid self-intersecting chords, and coloured according to the two-colour symmetry taxonomy:
\begin{itemize}
    \item \textbf{Symmetric Blocks (shaded blue)}: Polygons invariant under reflection on $\mathbb{Z}_{16}$, including regular squares $(4,4,4,4)$, rectangles $(1,7,1,7)$, and symmetric trapezoids.
    \item \textbf{Chiral Blocks (shaded pink)}: Polygons lacking reflectional symmetry on $\mathbb{Z}_{16}$, occurring in enantiomorphic pairs.
\end{itemize}

\subsection{The 5-Tier Natural Geometric Sorting Cascade}

To give the entire universe of $20{,}721$ tetramers a mathematically natural and visually striking progression, every cartouche is indexed according to a 5-tier hierarchical cascade:
\begin{enumerate}
    \item \textbf{Tier 1: Dihedral Orbit Size $|O|$}:
    Arranged in ascending order of dihedral orbit size:
    \[
    |O| \in \{4, 8, 16, 32\}
    \]
    The highest-symmetry orbits ($|O| = 4$ and $|O| = 8$) appear at the very head of the catalogue. Among all $20{,}721$ tetramers, exactly $2$ have orbit size $4$, $21$ have orbit size $8$, $622$ have orbit size $16$, and the remaining $20{,}076$ have orbit size $32$.
    
    \item \textbf{Tier 2: Symmetric Block Count $s$}:
    Within each orbit size, tetramers are sorted by the number of symmetric (achiral) blocks $s \in \{4, 3, 2, 1, 0\}$ in descending order ($4 \to 0$). This induces an unmistakable visual transition: pages begin in pure cobalt blue ($s = 4$), with pink chiral blocks appearing gradually until the final sections display completely chiral patterns ($s = 0$).
    
    \item \textbf{Tier 3: Chord Dispersion Energy $\mathcal{E}(P)$}:
    Sorted in descending order of circular chord dispersion energy:
    \[
    \mathcal{E}(P) = \sum_{B \in P} \sum_{x < y \in B} \min((y - x) \bmod 16, (x - y) \bmod 16)^2
    \]
    High energy corresponds to maximally regular, dispersed polygons with long chords; low energy corresponds to tightly clustered chords.
    
    \item \textbf{Tier 4: Canonical Stride Profile Multiset}:
    Grouping together tetramers whose inscribed quadrilaterals have identical unoriented cyclic stride profiles.
    
    \item \textbf{Tier 5: Canonical Lexicographic Tie-Breaker}:
    Canonical sorted block-set representation on $\mathbb{Z}_{16}$.
\end{enumerate}

\subsection{Comparative Publishing Scale Models}

Because printing all $20{,}721$ tetramers requires balancing visual resolution against total page count, three candidate publishing formats present themselves:
\begin{enumerate}
    \item \textbf{Model I: Micro-Atlas (220 Pages, 96 Tetramers / Page)}:
    Mosaic diameter $\varnothing = 7.0\text{ mm}$ ($r = 0.35\text{ cm}$). Laid out in a $4 \times 24$ cartouche grid (4 linear strips across the page, 24 rows down, giving $20{,}721 / 96 \approx 216$ pages). Dispensing with perimeter vertex dots preserves crisp legibility even at this compact scale.
    
    \item \textbf{Model II: Luxury Monograph (432 Pages, 48 Tetramers / Page)}:
    Mosaic diameter $\varnothing = 10.6\text{ mm}$ ($r = 0.53\text{ cm}$). Laid out in a $2 \times 24$ or $3 \times 16$ cartouche grid ($20{,}721 / 48 \approx 432$ pages). Strikes an ideal balance between individual mosaic legibility and library monograph format.
    
    \item \textbf{Model III: Gold-Tooled Leather-Bound (864 Pages in 2 Volumes, 24 Tetramers / Page)}:
    Mosaic diameter $\varnothing = 15.0\text{ mm}$ ($r = 0.75\text{ cm}$). Laid out in a $2 \times 12$ cartouche grid ($20{,}721 / 24 \approx 864$ pages). A grand exhibition format where every individual polygon edge and chord angle is immediately discernible.
\end{enumerate}

Figure~\ref{fig:tetramer-scale-comparison} displays the opening twelve tetramers $T_1 \dots T_{12}$ across these three publishing densities, with four linear cartouches demonstrated at each respective scale.

\begin{figure}[p]
\centering
\input{tkz/fig_tetramer_scale_comparison.tkz}
\caption{\label{fig:tetramer-scale-comparison}Comparative scale plate demonstrating the first twelve tetramers $T_1 \dots T_{12}$ across three candidate publishing models in horizontal linear cartouches ($1 \times 4$ strips): Tier~I ($T_1 \dots T_4$, Micro-Atlas scale, $\varnothing = 7.0\text{ mm}$); Tier~II ($T_5 \dots T_8$, Luxury Monograph scale, $\varnothing = 10.6\text{ mm}$); and Tier~III ($T_9 \dots T_{12}$, Gold-Tooled scale, $\varnothing = 15.0\text{ mm}$)}
\end{figure}
\clearpage

\subsection{Directory of Exhibited Sample Tetramers}

Table~\ref{tab:exhibited-tetramers} catalogues the twelve sample tetramers displayed in Figure~\ref{fig:tetramer-scale-comparison}, providing their algebraic invariants, stride profiles, and single-character hexadecimal partition-slash representations.

\begin{table}[htbp]
\centering
\footnotesize
\setlength{\tabcolsep}{3.5pt}
\begin{tabular}{rcccll}
\toprule
\textbf{ID} & \textbf{Model Scale} & \textbf{$|O|$} & \textbf{$\mathcal{E}$} & \textbf{Polygon Strides} & \textbf{Partition-Slash Representation} \\
\midrule
$T_1$  & Micro-Atlas   & $4$ & $720$ & $(1,6,3,6)^{\times 4}$ & \texttt{/017A/289F/36CD/45BE/} \\
$T_2$  & Micro-Atlas   & $4$ & $720$ & $(2,5,2,7)^{\times 4}$ & \texttt{/0279/18AF/35CE/46BD/} \\
$T_3$  & Micro-Atlas   & $8$ & $816$ & $(1,6,3,6)^{\times 2} + (1,7,1,7)^{\times 2}$ & \texttt{/017A/289F/34BC/56DE/} \\
$T_4$  & Micro-Atlas   & $8$ & $816$ & $(1,7,1,7)^{\times 2} + (2,5,2,7)^{\times 2}$ & \texttt{/0189/23AB/46DF/57CE/} \\
\midrule
$T_5$  & Luxury        & $8$ & $784$ & $(1,6,3,6)^{\times 2} + (1,7,1,7) + (3,5,3,5)$ & \texttt{/017A/289F/36BE/45CD/} \\
$T_6$  & Luxury        & $8$ & $784$ & $(1,7,1,7)^{\times 2} + (3,4,5,4)^{\times 2}$ & \texttt{/0189/26BF/37AE/45CD/} \\
$T_7$  & Luxury        & $8$ & $784$ & $(1,7,1,7) + (2,5,2,7)^{\times 2} + (3,5,3,5)$ & \texttt{/0189/27AF/35CE/46BD/} \\
$T_8$  & Luxury        & $8$ & $776$ & $(1,6,3,6)^{\times 2} + (2,6,2,6)^{\times 2}$ & \texttt{/017A/289F/35BD/46CE/} \\
\midrule
$T_9$  & Gold-Tooled   & $8$ & $776$ & $(2,5,2,7)^{\times 2} + (2,6,2,6)^{\times 2}$ & \texttt{/0279/18AF/35BD/46CE/} \\
$T_{10}$ & Gold-Tooled & $8$ & $736$ & $(2,5,4,5)^{\times 2} + (2,6,2,6) + (4,4,4,4)$ & \texttt{/027B/159D/38AF/46CE/} \\
$T_{11}$ & Gold-Tooled & $8$ & $720$ & $(1,4,7,4)^{\times 2} + (1,7,1,7)^{\times 2}$ & \texttt{/015C/23AB/489D/67EF/} \\
$T_{12}$ & Gold-Tooled & $8$ & $720$ & $(1,6,3,6)^{\times 2} + (2,5,2,7)^{\times 2}$ & \texttt{/017A/289F/35CE/46BD/} \\
\bottomrule
\end{tabular}
\caption{\label{tab:exhibited-tetramers}The twelve sample tetramers exhibited in Figure~\ref{fig:tetramer-scale-comparison}. All twelve patterns possess four symmetric blocks ($s = 4$). In the partition notation, vertices $0 \dots 15$ are represented as single hexadecimal characters \texttt{0} \dots \texttt{F}, indexing blocks between slashes}
\end{table}

The complete machine-readable index of all affine macro-orbits, formatted in this compact single-character notation with full cycle relabelings and chirality signatures, is preserved in the companion resource file \texttt{hexadecality\_index.txt}.


\appendix
% Appendix: The Small Universes: D_4 and D_9
% Complete Canonical Catalogues of Sister Universes

To appreciate the vastness of the $z_4^4$ universe on $\mathbb{Z}_{16}$, it is illuminating to examine its smaller sister universes: the $D_4$ universe of 2-block partitions on $\mathbb{Z}_4$ ($z_2^2$), and the $D_9$ universe of 3-block partitions on $\mathbb{Z}_9$ ($z_3^3$). Where $\mathbb{Z}_{16}$ generates tens of thousands of orbits across multiple pages, these sibling universes are completely and exhaustively surveyable in a single glance.

\section{The \texorpdfstring{$D_4$}{D4} Universe}
\label{sec:d4-universe}

The metric arena $\mathbb{Z}_4 = \{0, 1, 2, 3\}$ represents the square. A mosaic of shape $z_2^2$ is a set partition of the four vertices into two blocks of two vertices each.

\subsection{Dihedral Orbital Enumeration}

The total number of raw equal-block partitions of shape $z_2^2$ is:
\[
\frac{1}{2!}\binom{4}{2} = \frac{6}{2} = 3
\]
These three partitions are:
\begin{align*}
P_0 &= \{\{0, 2\}, \{1, 3\}\} \\
P_1 &= \{\{0, 1\}, \{2, 3\}\} \\
P_2 &= \{\{0, 3\}, \{1, 2\}\}
\end{align*}
The dihedral group $D_4$ of order $8$ (the symmetries of the square) acts transitively on the vertices. Under this action:
\begin{itemize}
    \item $P_0$ connects opposite vertices by orthogonal interior diagonals (stride profile $(2, 2)$). It is invariant under all four rotations and all four reflections of $D_4$. Its stabiliser has full order $|\operatorname{Stab}_{D_4}(P_0)| = 8$, yielding a singleton dihedral orbit of size $|O| = 1$.
    \item $P_1$ and $P_2$ connect adjacent boundary edges (stride profile $(1, 3)$). The rotation $r \colon x \mapsto x + 1 \pmod 4$ interchanges them ($r(P_1) = P_2$, $r(P_2) = P_1$). They merge into a single dihedral orbit $\{P_1, P_2\}$ of size $|O| = 2$, with stabiliser order $4$.
\end{itemize}
Hence, there are exactly two dihedral orbits under $D_4$.

\subsection{Quotient Collapse}

The unit group of the ring $\mathbb{Z}_4$ consists of the odd residues:
\[
(\mathbb{Z}_4)^\times = \{1, 3\} = \{1, -1\}
\]
Because reflection across vertex $0$ acts on $\mathbb{Z}_4$ as $x \mapsto -x \equiv 3x \pmod 4$, the inversion multiplier $\mu_3$ is already an element of the rigid dihedral reflection group: $\mu_3 \in D_4$.

Consequently, the outer affine automorphism group collapses completely:
\[
G_{\text{macro}} = \operatorname{Aff}(\mathbb{Z}_4) / D_4 \cong (\mathbb{Z}_4)^\times / \{\pm 1\} \cong C_1
\]
The macro-quotient is the trivial group of order $1$. Every dihedral orbit is therefore an affine macro-orbit of length $1$ (a monomer). There are no dimers or tetramers.

\begin{figure}[htbp]
\centering
\input{tkz/fig_d4_universe.tkz}
\caption{\label{fig:d4-universe}The complete $D_4$ universe: the two affine monomers of shape $z_2^2$ on $\mathbb{Z}_4$. $M_1$ is the regular orthogonal cross ($|O|=1$); $M_2$ represents the parallel edges ($|O|=2$)}
\end{figure}

The complete $D_4$ catalogue consists of precisely two monomers, illustrated in Figure~\ref{fig:d4-universe}:
\begin{enumerate}
    \item \textbf{Monomer $M_1$} ($|O|=1$): The regular coset partition of the ideal $2\mathbb{Z}_4$.
    \item \textbf{Monomer $M_2$} ($|O|=2$): The parallel chord pair.
\end{enumerate}
Machine-readable relational manifests linking both monomers to \texttt{polyatlas} compositions are tabulated in \path{data/z4_mosaics.tsv} and \path{data/z4_mosaic_blocks.tsv} (see Section~\ref{sec:hexdec-relational-datasets}).

\section{The \texorpdfstring{$D_9$}{D9} Universe}
\label{sec:d9-universe}

The metric arena $\mathbb{Z}_9 = \{0, 1, \dots, 8\}$ represents the regular nonagon. A mosaic of shape $z_3^3$ partitions the nine vertices into three blocks of three vertices each (three triangles).

\subsection{The 25 Dihedral Orbits}

The raw combinatorial count of equal-block partitions of shape $z_3^3$ is:
\[
\frac{1}{3!}\binom{9}{3, 3, 3} = \frac{1}{6} \cdot \frac{9!}{(3!)^3} = \frac{1680}{6} = 280
\]
The dihedral group $D_9$ has order $|D_9| = 18$ (nine rotations and nine reflections). Under Burnside's Lemma, the number of non-equivalent dihedral orbits is:
\[
\frac{1}{18} \sum_{g \in D_9} |\operatorname{Fix}(g)| = 25
\]
These $25$ dihedral orbits stratify into four stabiliser size classes:
\begin{itemize}
    \item \textbf{Size 1 (1 orbit)}: The unique coset partition of the ideal $3\mathbb{Z}_9$:
    \[
    P_0 = \{\{0, 3, 6\}, \{1, 4, 7\}, \{2, 5, 8\}\}
    \]
    Each block is a regular equilateral triangle with stride profile $(3, 3, 3)$ (shaded blue). Invariant under all $18$ symmetries of $D_9$.
    \item \textbf{Size 3 (3 orbits)}: Three orbits with order-$6$ stabiliser $\langle r^3, s \rangle$.
    \item \textbf{Size 9 (12 orbits)}: Twelve orbits with order-$2$ bilateral reflectional stabilisers.
    \item \textbf{Size 18 (9 orbits)}: Nine fully asymmetric chiral orbits ($|\operatorname{Stab}| = 1$).
\end{itemize}
The partition sum accounts for all $280$ configurations:
\[
1 \cdot 1 + 3 \cdot 3 + 12 \cdot 9 + 9 \cdot 18 = 1 + 9 + 108 + 162 = 280
\]

\subsection{The Outer Automorphism Group \texorpdfstring{$C_3$}{C3}}

The unit group $(\mathbb{Z}_9)^\times$ consists of the six integers coprime to $9$:
\[
(\mathbb{Z}_9)^\times = \{1, 2, 4, 5, 7, 8\} \cong C_6
\]
Because reflection across vertex $0$ acts as $x \mapsto -x \equiv 8x \pmod 9$, the subgroup $\{\pm 1\} = \{1, 8\}$ lies entirely within $D_9$. Taking the quotient yields the outer macro-automorphism group:
\[
G_{\text{macro}} = \operatorname{Aff}(\mathbb{Z}_9) / D_9 \cong (\mathbb{Z}_9)^\times / \{\pm 1\} \cong C_6 / C_2 \cong C_3
\]
This outer group $G_{\text{macro}} \cong C_3$ is generated by the dilation multiplier:
\[
\mu_2 \colon x \mapsto 2x \pmod 9
\]
with cycle structure $\mu_2 \colon 1 \mapsto 2 \mapsto 4 \mapsto 8 \equiv -1 \equiv 1 \pmod{G_{\text{macro}}}$.

Because $|G_{\text{macro}}| = 3$ is a prime number, the only possible divisors are $1$ and $3$. Thus, every macro-orbit in the $D_9$ universe is either:
\begin{itemize}
    \item A \textbf{Monomer} (fixed point of length $1$).
    \item A \textbf{Trimer} (a 3-cycle under $\mu_2 \colon O \mapsto \mu_2(O) \mapsto \mu_4(O) \mapsto O$).
\end{itemize}
Applying Burnside's Lemma to the action of $C_3$ on the $25$ dihedral orbits reveals that $\mu_2$ fixes only the regular coset partition $P_0$. The remaining $25 - 1 = 24$ dihedral orbits partition into exactly $24 / 3 = 8$ Trimers!

\subsection{The Complete Inventory of Trimers}

To maintain complete methodological uniformity across universes, the nonagonal trimers are catalogued under the exact same 5-tier natural sorting cascade established for the hexadecagonal tetramers in Section~\ref{sec:tetramers-and-scale-models}:
\begin{enumerate}
    \item \textbf{Tier 1: Dihedral Orbit Size $|O|$}: Ascending order $|O| \in \{3, 9, 18\}$.
    \item \textbf{Tier 2: Symmetric Block Count $s$}: Descending order $s \in \{3, 1, 0\}$, producing the characteristic visual transition from pure cobalt blue ($s=3$) through mixed chirality ($s=1$) to full chiral pink ($s=0$).
    \item \textbf{Tier 3: Circular Chord Dispersion Energy $\mathcal{E}$}: Descending order of total chord dispersion energy.
    \item \textbf{Tier 4: Canonical Stride Profile Multiset}: Grouping by unoriented cyclic stride profiles.
    \item \textbf{Tier 5: Lexicographic Tie-Breaker}: Sorting by canonical block elements.
\end{enumerate}
Furthermore, the orientation of each 3-strip is anchored by evaluating the 5-tier key across all three cycle phases, placing the phase of highest symmetry and dispersion energy at the leading $\mu_1$ position.

The complete census of eight trimers is catalogued in Table~\ref{tab:d9-trimers}. Together with the sole monomer $M_1$, these nine cartouches constitute the entire $D_9$ universe. The visual atlas is presented in Figure~\ref{fig:d9-universe}, where each trimer is displayed as a horizontal 3-strip tracking the Galois dilation sequence $\mu_1 \to \mu_2 \to \mu_4$. Full relational block manifests connecting all $25$ dihedral comosaics to \texttt{polyatlas} are tabulated in \path{data/z9_mosaics.tsv} and \path{data/z9_mosaic_blocks.tsv}.

\begin{table}[htbp]
\centering
\footnotesize
\setlength{\tabcolsep}{4.5pt}
\begin{tabular}{cccccc}
\toprule
\textbf{Trimer} & \textbf{$|O|$} & \textbf{$s$} & \textbf{$\mathcal{E}$} & \textbf{Canonical Strides} & \textbf{Galois Dilation Cycle ($\mu_1 \to \mu_2 \to \mu_4$)} \\
\midrule
$T_1$ & 3  & 3 & 99 & $(1, 4, 4)^3$ & \texttt{/015/267/348/} \texttt{[sss]} $\to$ \texttt{/012/345/678/} $\to$ \texttt{/024/168/357/} \\
$T_2$ & 9  & 3 & 90 & $(1, 4, 4)^2, (2, 2, 5)$ & \texttt{/015/237/468/} \texttt{[sss]} $\to$ \texttt{/012/378/456/} $\to$ \texttt{/024/138/567/} \\
$T_3$ & 9  & 1 & 91 & $(1, 4, 4), (2, 3, 4)^2$ & \texttt{/015/247/368/} \texttt{[snn]} $\to$ \texttt{/012/367/458/} $\to$ \texttt{/024/178/356/} \\
$T_4$ & 9  & 1 & 85 & $(1, 3, 5)^2, (1, 4, 4)$ & \texttt{/014/256/378/} \texttt{[nns]} $\to$ \texttt{/028/134/567/} $\to$ \texttt{/047/135/268/} \\
$T_5$ & 9  & 1 & 85 & $(2, 3, 4)^2, (3, 3, 3)$ & \texttt{/025/147/368/} \texttt{[nsn]} $\to$ \texttt{/014/258/367/} $\to$ \texttt{/028/147/356/} \\
$T_6$ & 18 & 1 & 88 & $(1, 3, 5), (1, 4, 4), (2, 3, 4)$ & \texttt{/014/267/358/} \texttt{[nsn]} $\to$ \texttt{/028/167/345/} $\to$ \texttt{/047/168/235/} \\
$T_7$ & 18 & 1 & 79 & $(1, 3, 5), (2, 2, 5), (2, 3, 4)$ & \texttt{/014/268/357/} \texttt{[nns]} $\to$ \texttt{/028/156/347/} $\to$ \texttt{/047/123/568/} \\
$T_8$ & 18 & 0 & 84 & $(1, 3, 5), (2, 3, 4)^2$ & \texttt{/014/257/368/} \texttt{[nnn]} $\to$ \texttt{/028/145/367/} $\to$ \texttt{/047/128/356/} \\
\bottomrule
\end{tabular}
\caption{\label{tab:d9-trimers}The complete catalogue of eight trimers of shape $z_3^3$ on $\mathbb{Z}_9$ ordered under the 5-tier natural sorting cascade ($|O| \uparrow$, $s \downarrow$, $\mathcal{E} \downarrow$, strides, lexicographical order)}
\end{table}

\begin{figure}[p]
\centering
\input{tkz/fig_d9_universe.tkz}
\caption{\label{fig:d9-universe}The complete $D_9$ universe: the sole Monomer ($M_1$, top) and the eight Trimers ($T_1 \dots T_8$, bottom), shown as horizontal 3-strips under the Galois generator $\mu_2 \colon x \mapsto 2x \pmod 9$}
\end{figure}

% Appendix B: Digital Resources and Computational Toolchain

This monograph is accompanied by an open suite of machine-readable datasets, relational polygon manifests, and interactive command-line utilities. All digital artifacts are designed to facilitate computational verification, algorithmic querying, and combinatorial research into cyclic lattice composition mosaics.

\section{The Atlas Index}
\label{sec:atlas-index}

The complete inventory of all $21{,}030$ affine macro-orbits is compiled in the machine-readable dataset:
\begin{center}
\path{hexdec/data/hexadecality_index.txt}
\end{center}
This comprehensive index contains $21{,}065$ lines covering the entire combinatorial universe of $83{,}488$ dihedral orbits ($14$ Monomers, $295$ Dimers, and $20{,}721$ Tetramers):

\begin{enumerate}
    \item \textbf{Record Identifier}:
    Prefixed by orbit class and index: \texttt{M01}--\texttt{M14} for monomers, \texttt{D001}--\texttt{D295} for dimers, and \texttt{T00001}--\texttt{T20721} for tetramers.
    
    \item \textbf{Physical Metric Invariants}:
    Each record tabulates the dihedral orbit size $|O| \in \{1, 4, 16, 32\}$, the block bilateral symmetry count $s \in \{0, \dots, 4\}$, and the total chord energy $E$.
    
    \item \textbf{Partition-Slash Notation}:
    Partitions are rendered in compact hexadecimal slash notation, where vertices $10 \dots 15$ are represented by uppercase letters $\text{A} \dots \text{F}$. For example:
    \begin{verbatim}
    /0279/18AF/35CE/46BD/
    \end{verbatim}
    denotes the partition $\{\{0, 2, 7, 9\}, \{1, 8, 10, 15\}, \{3, 5, 12, 14\}, \{4, 6, 11, 13\}\}$.
    
    \item \textbf{Block Chirality Signatures}:
    Appended to each partition is a 4-letter signature in brackets, such as \texttt{[ssss]}, \texttt{[ssnn]}, or \texttt{[nnnn]}. Each character corresponds to one block:
    \begin{itemize}
        \item \texttt{s} (symmetric): the block possesses an internal bilateral reflection axis in $\mathbb{Z}_{16}$, rendering it achiral and non-orientable (shaded blue in plates).
        \item \texttt{n} (non-symmetric): the block lacks any internal reflection axis, rendering it intrinsically chiral (shaded pink in plates).
    \end{itemize}
    
    \item \textbf{Galois Dilation Cycle Transitions}:
    For dimers and tetramers, the line tracks the transition cycle under the affine generator $\mu_3 \colon x \mapsto 3x \pmod{16}$:
    \begin{verbatim}
    D00001 Dimer |O|=4 s=4 E=640 /028A/139B/46CE/57DF/ [ssss]
      -> /068E/179F/24AC/35BD/ [ssss]
    \end{verbatim}
\end{enumerate}

\section{The Mosaical Orbit Resolver}
\label{sec:lookup-tool}

The command-line tool \texttt{hexdec/tools/lookup\_mosaic.py} provides an automated resolution engine for any user-supplied set partition of $\mathbb{Z}_{16}$. 

Given an arbitrary partition string in partition-slash format, the tool:
\begin{enumerate}
    \item Verifies that the string forms a valid, non-overlapping partition of $\{0, \dots, 15\}$ into four blocks of size four.
    \item Computes the full $32$-element dihedral orbit under $D_{16}$ and extracts the lexicographical canonical representative.
    \item Identifies the stabiliser subgroup and determines the dihedral orbit size $|O| \in \{1, 4, 16, 32\}$.
    \item Calculates the chord energy $E$ and the block chirality signature.
    \item Resolves the parent macro-orbit (Monomer, Dimer, or Tetramer) and its primary catalogue index.
    \item Traces the full macro-orbit under Galois dilation $\mu_3 \colon x \mapsto 3x \pmod{16}$, printing both the geometrically dilated chord diagrams and the corresponding canonical dihedral indices.
\end{enumerate}

\begin{example}
Querying the tetramer \texttt{/0279/18AF/35CE/46BD/} yields:
{\small
\begin{verbatim}
$ lookup-mosaic "/0279/18AF/35CE/46BD/"
======================================================================
HEXADECALITY: COMPOSITION MOSAIC ORBIT & CARTOUCHE RESOLVER
======================================================================
Input Partition:       /0279/18AF/35CE/46BD/  [ssss]
Canonical Dihedral D16:/0279/18AF/35CE/46BD/  [ssss]
Dihedral Orbit Size:   |O| = 4 (Stabiliser order: 8)
Symmetric Blocks:      s = 4 of 4 (Achiral)
Chord Energy:          E = 720
----------------------------------------------------------------------
Affine Macro-Orbit Type: Tetramer (4-Cycle)

Orbital Comosaics Cycle (Physical Dilation vs Canonical Rep):
  Phase 1 (mu_1 = id) : /0279/18AF/35CE/46BD/ [ssss] (Canonical)
  Phase 2 (mu_3 = x3) : /056B/127C/38DE/49AF/ [ssss]
                        (Canonical: /016B/27CD/389E/45AF/)
  Phase 3 (mu_7 = x7) : /01EF/2345/6789/ABCD/ [ssss]
                        (Canonical: /0123/4567/89AB/CDEF/)
  Phase 4 (mu_5 = x5) : /03AD/147E/258B/69CF/ [ssss]
                        (Canonical: /0369/18BE/25CF/47AD/)
======================================================================
\end{verbatim}}
\end{example}

\section{The Relational Polygon Datasets and Embedding Manifests}
\label{sec:hexdec-relational-datasets}

To establish a bidirectional bridge between composition mosaics and the canonical inscribed polygon census of \texttt{polyatlas}, the entire combinatorial corpus across all three metric universes ($\mathbb{Z}_4$, $\mathbb{Z}_9$, and $\mathbb{Z}_{16}$) is provided as normalised, machine-readable tab-separated resource files housed in \path{hexdec/data/}.

\subsection{Relational Architecture and Syntax}

Each universe is structured as a relational pair:
\begin{enumerate}
    \item \textbf{Mosaic-Level Catalogue (\path{data/z*_mosaics.tsv})}:
    Tabulates the global dihedral and affine invariants ($|O|$, $|\operatorname{Stab}|$, $s$, $E$, chirality signature), the constituent polygon multiset profile, and the collective \emph{block manifest}.
    
    \item \textbf{Block-Level Relational Bridge (\path{data/z*_mosaic_blocks.tsv})}:
    Links each individual block in every canonical mosaic directly to its parent inscribed polygon class in \path{polyatlas/data/inngons.tsv}.
\end{enumerate}

Every block embedding is encoded by the standardised manifest descriptor:
\begin{equation}
\label{eqn:hexdec-block-manifest}
\text{[Class]}@\text{[Root]}[\text{Chirality}]:\text{[Realised Strides]}
\end{equation}
where \textbf{Class} denotes the parent polygon class $k\text{-}i$ defined in \texttt{polyatlas}; \textbf{Root} is the starting vertex $v_0 = \min(B_j)$ around the clock face; \textbf{Chirality} specifies the orientation sign ($+$ for direct rotational alignment, $-$ for retrograde reflection); and \textbf{Realised Strides} lists the forward perimeter strides from $v_0$ along the block.

\subsection{Multi-Universe Relational Directory}

The companion datasets span three cyclic universes:
\begin{itemize}
    \item \textbf{$D_4$ Universe ($z_2^2$ on $\mathbb{Z}_4$)}:
    \path{data/z4_mosaics.tsv} ($2$ orbits) and \path{data/z4_mosaic_blocks.tsv} ($4$ block records). Both digon classes ($2\text{-}1$ and $2\text{-}2$) are realised.
    
    \item \textbf{$D_9$ Universe ($z_3^3$ on $\mathbb{Z}_9$)}:
    \path{data/z9_mosaics.tsv} ($25$ dihedral comosaics across $9$ macro-orbits) and \path{data/z9_mosaic_blocks.tsv} ($75$ block records). All $7$ triangle classes ($3\text{-}1$ through $3\text{-}7$) are fully realised.
    
    \item \textbf{$D_{16}$ Universe ($z_4^4$ on $\mathbb{Z}_{16}$)}:
    \begin{itemize}
        \item \textbf{Primary Macro-Orbits}:
        \path{data/z16_macro_mosaics.tsv} ($21{,}030$ records) and\\
        \path{data/z16_macro_mosaic_blocks.tsv} ($84{,}120$ block records).
        \item \textbf{Exhaustive Dihedral Comosaics}:
        \path{data/z16_mosaics.tsv} ($83{,}488$ records) and\\
        \path{data/z16_mosaic_blocks.tsv} ($333{,}952$ block records).
    \end{itemize}
\end{itemize}

\subsection{The Injection Property and Cascade Filter}

Unlike unconstrained partitions, the mosaics treated here are equal-block partitions ($z_k^m$), representing structured injections into the general compositional universe:
\begin{itemize}
    \item \textbf{Enantiomorphic Realisation}: In \texttt{polyatlas}, chiral polygons appear strictly in their canonical direct orientation. In mosaics, chiral blocks appear in both direct ($+$) and retrograde ($-$) configurations (e.g.\ class $4\text{-}68$ exhibits $3{,}303$ direct and $2{,}563$ retrograde blocks across $\mathbb{Z}_{16}$).
    \item \textbf{Universal Surjectivity}: Across the complete $83{,}488$ comosaics of $\mathbb{Z}_{16}$, all $72$ tetragon classes from \texttt{polyatlas} are realised ($100\%$ coverage).
    \item \textbf{Sorting Cascade Dispersion}: In the $21{,}030$ primary macro-orbit representatives, exactly $70$ of the $72$ classes appear. Classes $4\text{-}1$ ($/0123/$, strides $111\text{D}$) and $4\text{-}3$ ($/0134/$, strides $121\text{C}$) are filtered out of the leading phase because the 5-tier natural sorting cascade prioritises high chord dispersion energy $E$, systematically shifting clustered blocks into dilated comosaic phases ($\mu_3$, $\mu_7$, or $\mu_5$).
\end{itemize}

\begin{table}[htbp]
\centering
\caption{\label{tab:hexdec-sample-manifests}Exemplar Composition Mosaics and their Relational Manifests}
\resizebox{\linewidth}{!}{%
\begin{tabular}{l c c c l l}
\toprule
Identifier & Type & $|O|$ & $s$ & Polygon Profile & Relational Embedding Manifest \\
\midrule
$M_1$ ($\mathbb{Z}_4$) & Monomer & 1 & 2 & $2 \times 2\text{-}2$ & \path{2-2@0+:22/2-2@1+:22} \\
$M_2$ ($\mathbb{Z}_4$) & Monomer & 2 & 2 & $2 \times 2\text{-}1$ & \path{2-1@0+:13/2-1@2+:13} \\
\midrule
$M_1$ ($\mathbb{Z}_9$) & Monomer & 1 & 3 & $3 \times 3\text{-}4$ & \path{3-4@0+:333/3-4@1+:333/3-4@2+:333} \\
$T_1$ ($\mathbb{Z}_9$) & Trimer & 3 & 3 & $3 \times 3\text{-}2$ & \path{3-2@0+:144/3-2@2+:414/3-2@3+:144} \\
$T_3$ ($\mathbb{Z}_9$) & Trimer & 9 & 1 & $3\text{-}2 + 2 \times 3\text{-}7$ & \path{3-2@0+:144/3-7@2+:234/3-7@3-:324} \\
\midrule
$M_{01}$ ($\mathbb{Z}_{16}$) & Monomer & 1 & 4 & $4 \times 4\text{-}28$ & \path{4-28@0+:4444/4-28@1+:4444/4-28@2+:4444/4-28@3+:4444} \\
$M_{02}$ ($\mathbb{Z}_{16}$) & Monomer & 4 & 4 & $2 \times 4\text{-}13 + 2 \times 4\text{-}27$ & \path{4-13@0+:1717/4-27@2+:5353/4-27@3+:3535/4-13@4+:1717} \\
$M_{05}$ ($\mathbb{Z}_{16}$) & Monomer & 16 & 4 & $4\text{-}13+4\text{-}22+4\text{-}27+4\text{-}28$ & \path{4-13@0+:1717/4-27@2+:3535/4-28@3+:4444/4-22@4+:2626} \\
$T_{00001}$ ($\mathbb{Z}_{16}$) & Tetramer & 4 & 4 & $4 \times 4\text{-}12$ & \path{4-12@0+:1636/4-12@2+:1636/4-12@3+:3616/4-12@4+:1636} \\
\bottomrule
\end{tabular}%
}
\end{table}

% Appendix C: A Polyrhythmic Percussion Suite

The geometric partition of the sixteen vertices $\mathbb{Z}_{16} = \{0, 1, \dots, 15\}$ admits an immediate acoustic realisation as a $16$-step rhythmic measure (four beats of sixteenth notes in common $4/4$ time).

Because each mosaic $P = \{B_0, B_1, B_2, B_3\}$ forms a partition of the metric grid, at each sixteenth-note step $t \in \{0, \dots, 15\}$ exactly one block contains $t$. When each block is assigned a distinct percussive instrument, the four voices never sound simultaneously. They generate a dense, interlocking hocket: an uninterrupted composite stream of sixteenth notes where melodic and rhythmic contours emerge purely from timbre.

\section{Acoustic Mapping and Voice Assignment}
\label{sec:percussion-suite}

Fourteen Standard MIDI files (\texttt{monomer1.mid} through \texttt{monomere.mid}, indexed in hexadecimal $1 \dots \text{e}$) form a complete percussion suite synthesised on General MIDI Channel~10:
\begin{itemize}
    \item \textbf{Voice 0} (Block containing downbeat vertex $0$): Bass Drum 1 (Kick, MIDI note 36), anchoring the measure's pulse.
    \item \textbf{Voice 1}: Acoustic Snare (MIDI note 38), providing the mid-frequency backbeat.
    \item \textbf{Voice 2}: Closed Hi-Hat (MIDI note 42), supplying high-frequency metallic articulation.
    \item \textbf{Voice 3}: Claves (MIDI note 75), cutting through with a sharp, resonant wooden click.
\end{itemize}

Each monomer exhibits a distinct metric character governed by its dihedral and affine stabilisers:
\begin{itemize}
    \item \textbf{Regular Pulse ($M_1$)}:
    Because $M_1$ consists of the four cosets of $4\mathbb{Z}_{16}$, each instrument plays an isochronous stream of four quarter notes, offset by single sixteenth notes. It sounds as four interwoven "four-on-the-floor" pulses.
    
    \item \textbf{Palindromic Grooves ($M_2, M_3, M_4$)}:
    Owing to their order-$8$ dihedral stabilisers, these rhythms display exact half-measure translational symmetry ($t \mapsto t + 8$) and bilateral retrogrades. The kick and snare form balanced, syncopated funk phrases with predictable metric counterpoint.
    
    \item \textbf{Reflectional Syncopations ($M_5, M_6, M_7, M_8$)}:
    $M_8$ exhibits an inversion groove ($t \mapsto t + 8$), swapping instrument pairings halfway through the bar, while $M_5, M_6, M_7$ possess metric retrograde axes where the sixteenth-note pattern reads identically forwards and backwards around specific beats.
    
    \item \textbf{Chiral Breakbeats ($M_9 \dots M_{14}$)}:
    Having no Euclidean symmetries, these rhythms sound completely irregular and unpredictable in linear time. Their hidden periodicity only reveals itself under affine multiplier decimation: sub-sampling the rhythm at every third sixteenth note ($\mu_3 \colon t \mapsto 3t \pmod{16}$) recovers the exact original rhythmic profile.
\end{itemize}

\section{The Companion Synthesizer (\texttt{midipartner})}
\label{sec:percussion-synthesizer}

The percussion synthesiser \texttt{midipartner} (\path{tools/midipartner.py}, aliased as \texttt{midipart}) synthesises any valid mosaic partition into a multi-track General MIDI percussion loop on Channel~10.

\subsection{Catalogue Specifiers}
For hexadecimological study, \texttt{midipartner} directly parses canonical catalogue specifiers for all $21{,}030$ affine macro-orbits ($83{,}488$ dihedral orbits) tabulated in the atlas:
\begin{itemize}
    \item \textbf{Tetramer Patterns}: The specifier \texttt{T01152.3} selects pattern~3 of tetramer \texttt{T01152} under $\mu_3$ dilation. Omitting the pattern suffix (\texttt{T01152}) defaults to the primary representative (Phase~1):
{\small
\begin{verbatim}
midipartner T01152.3 --output T01152_3.mid --bpm 63 --bars 8
midipartner T01152 --bpm 63 --bars 8
\end{verbatim}}
    \item \textbf{Dimer Patterns}: The specifier \texttt{D001.2} selects the conjugate comosaic of dimer \texttt{D001} (omitting the suffix defaults to Phase~1):
{\small
\begin{verbatim}
midipartner D001.2 --bpm 90 --bars 4
\end{verbatim}}
    \item \textbf{Monomers}: All 14 monomers are addressed directly by record:
{\small
\begin{verbatim}
midipartner M01 --bpm 110 --bars 4
\end{verbatim}}
\end{itemize}
The utility scans the precomputed index in \path{hexdec/data/hexadecality_index.txt}, extracting the exact partition and chirality signature instantaneously without requiring runtime combinatorial generation.

\subsection{Arbitrary Partitions and Metric Grids}
Beyond catalogue lookups on $\mathbb{Z}_{16}$, \texttt{midipartner} provides general set partition parsing:
\begin{itemize}
    \item \textbf{Explicit Slash Notation}: Partitions may be supplied directly as positional arguments or via \texttt{-p}/\texttt{--partition}:
{\small
\begin{verbatim}
midipartner "/0149/26AE/37BF/58CD/" -o chiral_breakbeat.mid
\end{verbatim}}
    \item \textbf{Generalised Metric Grids and Odd Time Signatures ($n \le 36$)}:
    Partitions on cyclic groups $\mathbb{Z}_n$ up to $n=36$ are encoded using Base-36 alphanumeric characters (\texttt{0}--\texttt{9} for $0 \dots 9$ and \texttt{A}--\texttt{Z} / \texttt{a}--\texttt{z} for $10 \dots 35$). Rather than forcing odd cycles into an artificial common-time ($4/4$) grid, the synthesiser embeds the exact intrinsic **$n/16$ time signature** with semiquaver metronome resolution. For example, an $11$-step partition on $\mathbb{Z}_{11}$ is signatured as $11/16$, a nonagonal trimer on $\mathbb{Z}_9$ as $9/16$, and a heptagonal rhythm on $\mathbb{Z}_7$ as $7/16$:
{\small
\begin{verbatim}
midipartner "/0123/4567/89A/" -o odd_11_16.mid
midipartner "/012/345/678/" --bpm 108 -o trimer_z9.mid
midipartner "/012345/6789AB/CDEFGH/IJKLMN/OPQRST/UVWXYZ/" \
  -o z36.mid
\end{verbatim}}
    Metric grids with $n > 36$ are supported using hyphen- or space-separated multi-digit integers (e.g.\ \texttt{/0-1-18-36/2-3-19-37/.../}).

    \item \textbf{Additive Metre Groupings and Dynamic Accents}:
    For odd metres, \texttt{midipartner} supports custom metric pulse groupings via the \texttt{--grouping} flag (or explicit accents via \texttt{--accents}). For instance, an $11/16$ measure can be articulated with a $3+3+3+2$ or $4+4+3$ additive pulse:
{\small
\begin{verbatim}
midipartner "/0123/4567/89A/" --grouping 3,3,3,2 -o balkan_11_16.mid
midipartner "/0123/4567/89A/" --time-sig 11/16 --accents 0,3,6,9
\end{verbatim}}
    This places secondary acoustic accents at the sub-beat boundaries while preserving the primary measure downbeat on step~0, preventing unintended quarter-note common-time pulsations.
\end{itemize}

\subsection{Instrument Nomination}
To facilitate subtractive acoustic analysis, solo transcription, or isolating specific polyrhythmic strata, \texttt{midipartner} allows nominating a "silent" instrument for any partition part from Part~1 to Part~$k$:
\begin{itemize}
    \item \textbf{Dedicated Mute Flag}: Passing \texttt{--silent} (or \texttt{--mute}) with 1-indexed part numbers to suppress selected voices (e.g.\ \texttt{--silent 1} to mute the downbeat kick and audition the accompaniment, or \texttt{--silent 1 3}):
{\small
\begin{verbatim}
midipartner T01152.3 --silent 1 -o T01152_3_no_kick.mid
\end{verbatim}}
    \item \textbf{Instrument Palette Assignment}: Specifying \texttt{silent}, \texttt{rest}, \texttt{none}, or \texttt{0} via the \texttt{-i} palette flag:
{\small
\begin{verbatim}
midipartner T01152.3 -i silent 38 42 75 -o p1_muted.mid
midipartner T01152.3 -i silent snare hihat claves
\end{verbatim}}
\end{itemize}
Muted voices emit no MIDI note events while preserving the exact metric bar length of the ensemble loop.

\section{Score Engraving with \texttt{xmusicpartner}}
\label{sec:percussion-musicxml}

While Standard MIDI files provide instantaneous multi-track acoustic playback, they possess inherent structural limitations for music notation. Standard MIDI records only asynchronous timestamped note-on and note-off events; it contains no concept of musical rests, metric beam groupings, unpitched percussion staves, or score typography. When raw MIDI files are imported into notation software, heuristic quantisation algorithms frequently misinterpret silent intervals between hockets as overlapping tied notes, while failing to beam odd or additive metres ($7/16, 9/16, 11/16, 13/16$) correctly.

To provide a fully vendor-neutral notation pipeline, the companion suite includes the score generator \texttt{xmusicpartner} (housed at \path{tools/xmusicpartner.py}, aliased as \texttt{xmusicpart}). The utility targets the open \textbf{MusicXML 3.1/4.0} standard maintained by the W3C Music Notation Community Group. Generated scores import deterministically into any standard-compliant notation software, scorewriter, or musicological framework (including open-source engravers such as MuseScore and GNU LilyPond via \texttt{musicxml2ly}, commercial environments such as Steinberg Dorico and Avid Sibelius, and computational musicology toolkits such as \texttt{music21}).

\subsection{Notational Semantics and Architecture}

The exporter bridges cyclic combinatorial geometry and score engraving through several core mechanisms:
\begin{enumerate}
    \item \textbf{Deterministic Notes and Rests}:
    Each sixteenth-note time step $t \in \{0, \dots, n-1\}$ is evaluated across the partition. Active voices emit explicit unpitched semiquaver notes, while idle voices emit explicit sixteenth rests (\texttt{<rest/><type>16th</type>}). This eliminates quantisation guesswork and guarantees faithful score engraving.

    \item \textbf{Additive Metres and Metric Beam Trees}:
    For any cyclic modulus $n$, the \texttt{--grouping} option constructs explicit hierarchical beam trees (\texttt{begin}, \texttt{continue}, \texttt{end}). An $11/16$ metre specified with \texttt{--grouping 3,3,3,2} is automatically beamed into distinct metric clusters of $3$, $3$, $3$, and $2$ semiquavers, rather than collapsing into arbitrary four-note blocks or fragmented flags.

    \item \textbf{Percussion Clefs and Display Positions}:
    Measures are initialised with standard W3C unpitched percussion clefs (\texttt{<clef><sign>percussion</sign></clef>}), mapping the four mosaic voices to standard percussion staff positions: Bass Drum 1 on Space~1 ($F_4$), Acoustic Snare on Space~3 ($C_5$), Closed Hi-Hat with cross noteheads above Line~5 ($G_5$), and Claves on Space~4 ($E_5$).

    \item \textbf{Score Topologies (\texttt{--score-type})}:
    The generator provides two distinct structural score representations:
    \begin{itemize}
        \item \textbf{Multi-Staff Ensemble (\texttt{--score-type multi}, default)}: Assigns each polygon block to an independent percussion staff, clearly exhibiting the individual polyrhythmic stride and voice shape ($\mathbf{z}_k$) of each part.
        \item \textbf{Consolidated Drum Kit (\texttt{--score-type kit})}: Collates all interlocking voices into a single standard 5-line drum set staff, visually displaying the composite, gapless hocket stream.
    \end{itemize}

    \item \textbf{Combinatorial and Catalogue Metadata}:
    Every generated score embeds provenance in its XML header, recording the cyclic universe $\mathbb{Z}_n$, macro-orbit catalogue identifier (e.g.\ \texttt{T01152}, \texttt{D001}, \texttt{M01}), dilation phase index, partition shape, and composer credit.
\end{enumerate}

\subsection{Command-Line Execution}

The tool accepts both canonical atlas specifiers and arbitrary set partition strings:
\begin{itemize}
    \item \textbf{Atlas Catalogue Queries}:
{\small
\begin{verbatim}
xmusicpartner T01152.3 --bars 2 -o T01152_3.musicxml
xmusicpartner M01 --bars 4 -o monomer1.musicxml
xmusicpartner D001.2 --score-type kit -o dimer1_2_kit.musicxml
\end{verbatim}}

    \item \textbf{Additive Odd-Metre Partitions}:
{\small
\begin{verbatim}
xmusicpartner "/0123/4567/89A/" --grouping 3,3,3,2 -o odd11.musicxml
xmusicpartner "/0149/26AE/37BF/58CD/" --title "Chiral Breakbeat" \
  --score-type kit -o breakbeat.musicxml
\end{verbatim}}
\end{itemize}


\FloatBarrier
\bibliographystyle{plain}
\bibliography{gemai}

\end{document}
